% !TeX root = ../概率论与数理统计笔记.tex
% 由主文档通过 \input 载入；请编译主文档。
\setcounter{chapter}{18} % 保留原书章号。
\chapter{随机过程初步（选读）}
\label{chap:ps-processes}

本章把单个随机变量或有限随机向量推广为随时间变化的随机系统。
先修内容为联合分布、条件期望与极限定理；学习目标是掌握有限状态马尔可夫链和齐次泊松过程的基本模型。

\section{随机过程与时间结构}
\label{sec:ps-process-language}
定义\emph{随机过程}（stochastic process）、指标集、状态空间、有限维分布与样本路径。
区分离散时间和连续时间、离散状态和连续状态，并通过随机游走说明时间相关性。
只给出均值和协方差通常不能确定整个过程的分布。

\section{有限状态马尔可夫链}
\label{sec:ps-markov-chains}
定义\emph{马尔可夫性质}（Markov property）与时间齐次转移矩阵，说明初始分布也是模型的一部分。
推导多步转移概率与查普曼--柯尔莫哥洛夫方程（Chapman--Kolmogorov equations）。
介绍可达、互通、闭类、吸收状态、常返与暂留以及周期的含义，主线限定在有限状态空间。

\section{平稳分布、极限分布与长期平均}
\label{sec:ps-stationarity}
定义\emph{平稳分布}（stationary distribution）并求解归一化线性方程。
说明有限不可约链具有唯一平稳分布；从任意初始分布收敛到它还需非周期性。
用周期链区分平稳分布与边缘分布极限，并说明有限不可约链的长期时间平均结论不要求非周期性。
介绍细致平衡与可逆性，为第\ref{chap:ps-computation}章的模拟方法作准备。

\section{泊松过程与等待时间}
\label{sec:ps-poisson-process}
定义强度为 $\lambda>0$ 的\emph{齐次泊松过程}（homogeneous Poisson process），说明起点、独立增量、平稳增量及计数分布。
研究计数与到达时刻的联系，推导独立指数间隔和伽马到达时间分布。
讨论独立泊松过程叠加及独立标记下的稀疏化，明确这些独立性条件不可省略。

\section{例题、习题与学习检查}
规划天气状态的简化链、带吸收边界的随机游走、泊松到达的等待与计数问题。
习题要求写明初始分布和转移规则、检查平稳方程，并构造平稳分布存在但分布不收敛的周期例子。
鞅与布朗运动仅列为后续课程方向，不作为本书主干的先修要求。






